\documentclass[paper=a4,twoside=false,BCOR=0mm,DIV=calc,fontsize=12pt]{scrartcl}

\usepackage[automark,headsepline]{scrpage2}
% \usepackage{xunicode,fontspec,xltxtra}
\usepackage[english,german]{babel}

\usepackage[T1]{fontenc}
\usepackage[utf8x]{inputenc}

\usepackage{graphicx}
\usepackage{lastpage}
\usepackage{listings}

%Schöne Tabellen
\usepackage{booktabs}

%\usepackage{longtable}

\usepackage{multirow}

% Kein einrücken und Abstand zwischen Absätzen.
\usepackage{parskip}
\setlength{\parindent}{0cm}

% Hyperlinks und url's
\usepackage[hidelinks]{hyperref}



% \usepackage[debug]{libertine}
% \setromanfont[Mapping=tex-text]{Linux Libertine O}
% \setsansfont[Mapping=tex-text]{Linux Biolinum O}
% \setmonofont[Mapping=tex-text]{Liberation Mono}

\pagestyle{scrheadings}

\clearscrheadfoot
\ihead{\headmark}
\ohead{Seite\pagemark\ von \pageref{LastPage}}
\ifoot{Business frameworks}
\cfoot{{\includegraphics[width=1.5cm]{./img/CC3.png}}}
\ofoot{\today}

\renewcommand*{\pnumfont}{
	\normalfont\rmfamily\slshape
}

\KOMAoptions{draft=false}
\KOMAoptions{DIV=last}

\begin{document}

% --- Titelseite --- %
\begin{titlepage}
	\enlargethispage{3cm}
	\begin{raggedright}
	\begin{picture}(0,0)
		\put(-30,14){\includegraphics[width=7cm]{./img/fhnw-technik-head}}
	\end{picture}

	\vspace*{6cm}
	{\Huge\bfseries\sf
		MAS-IT\\[1.7ex]
	}
	{\Large\bfseries\sf
		Business frameworks\\[2.2ex]
	}
	{\large\bfseries\sf
		Masterthesis von\\[1.5ex]
		Etienne Rebetez\\[1.5ex]
	}
	\vspace*{1.5cm}
	{\large\bfseries\sf
		FHNW\\[1.5ex]
		Hochschule für Technik\\[1.5ex]
		Studiengang MAS-IT\\[1.5ex]
		Betreuender Dozent:\\[1.5ex]
		Dr. sc. techn. Ronald Tanner\\[1.5ex]
	}
	\vspace*{2cm}
	{\large\bfseries\sf
		Bern, \today\\
	}
	\end{raggedright}
\end{titlepage}

\newpage
% --- Zusammenfassung --- %
\section*{Summary}
Text

% --- Vorwort --- %
\section*{Vorwort}
Text
% --- Inhaltsverzeichnis --- %
\newpage
	\tableofcontents

% --- Einleitung --- %
\newpage
\section{Einleitung}
Laborprozesse. Was versteht man eigentlich darunter? 
Der stark abstrahierte Fall ist der folgende. Ein Physische Probe kommt in ein Analyselabor, bestimmte physikalische Parameter werden 
von der Probe erhoben und diese Resultate werden dem Auftraggeber bekanntgegeben.

In diesem Bericht wird ein Labor als Qualitätsicherungslabor betrachtet. Das heisst, die Prozesse und Methoden sind vorgeben und können bzw. dürfen nicht 
einfach abgeändert werden.

Traditionell wird die Probenverwaltung von einem LIMS (Labor Informations und Management System) geregelt. Traditionellerweise ist es auch so, dass
sobald die Probe welche analysiert werden muss im Labor steht, viele Schritte Manuell geschehen. Dies ist aufwendig und Fehleranfällig. 
Es darf natürlich nicht vergessen werden, dass unter Manuell auch Handarbeiten wie einwägen, Lösungen erstellen oder Analysegeräte befüllen fallen.
Für diese Aufgaben wird es noch lange Menschen geben welche diese Arbeit ausführen müssen.

Will man nun diese Laborprozesse Automatisieren, geht es vor allem um die Datenübertragung. 
Was muss gemessen werden, welche Methode soll verwendet werden,
welches qualifizierte Gerät darf verwendet werden, 
wie werden die Daten verrechnet und schlussendlich wie werden die Resultate wieder an das nächst höhere System weitergegeben.
Gleichzeitig solle das System den Laborarbeiter unterstützen und die durchgeführten Schritte sollten in alle nachvollziehbar sein.


<Geräte Software>



Menschen sind ausserordentlich gut darin Informationen 


\section{Die Laborumgebung}
Im Labor gibt es nun verschiedenste Methoden welche unterschiedliche Historische Ursprünge haben. So ist eine Chemische Titration ein
vollkommen anderes Konzept als z.B eine Biochemische Elisa Analyse.
Dies ergibt ein äusserst Heterogene Landschaft an Analysegeräten mit und ohne Software. Resultate oder Rohdaten liegen in unterschiedlichsten Formaten und Layouts vor.





\section{Konzept}
Um eine Automatisierung in diesem Heterogenen Umfeld realisieren zu können, braucht es eine Engine die einerseits den Status der einzelnen
Abläufe kennt und auch weiss welcher Businessprozess als nächster aufgerufen werden soll. Diese Statemachien oder Sheduler währe somit das
Herzstück der Automatisierung. Häufig fällt auch der Begriff SOA (Service Oriented Architecture).

Das Labor ist lange nicht der einzige Bereich welcher eine Automatisierung verlangt. Andere Branchen habe grundsätzlich das selbe
Bedürfnis. Daher gibt es schon eine reihe von Produkten welche diese Statemachienen bereits anbieten.

Die Konfiguration dieser Engines geschieht über standardisierte XML Files. Dabei haben sich zwei Beschreibungssprachen etabliert:
\begin{itemize}
 \item BPEL
 \item BPMN
\end{itemize}

Beide sind offene Standards vom OMG \cite{omg}.

Neben dieser Zentralen Steuerungseinheit, soll es dann auch möglich sein, Komponenten wie Webapplikationen (Formulare, Listen, etc),
Reports oder Mobile Geräte einbinden zu können. In den Labors ist es zudem Historisch so, das viele MS Excel Berechnungen existieren. Diese
sollten, dann auch in den Prozessen eingebunden werden können.


\section{BPEL}
Die BPEL (Business Process Execution Language) Beschreibungsprache kommt aus der Informatikwelt. Diese wird hat Beschreibungen zu
Schleifen, Bedingugnen und anderes Konstrukte die man aus der Programmierung kennt. 


\section{BPMN}
BPMN (Business Process Model and Notation) kommt aus dem Business und war gedacht erstmals auf dem Papier die Prozesse einheitlich
beschreiben zu können. Grob gesagt, handelt es sich um Flowcharts. 
Mit der Version 2 vom BPMN können diese Flowcharts nun auch programmatisch ausgeführt werden. 




\section{Engines}
Die Engine ist nun eine Implementierung einer solchen Business Prozess Beschreibung. Die Engines erfüllen nun die Logik. 
\begin{itemize}
  \item Welche Prozesse gibt es?
  \item Welche sind aktiv und in welchem zustand?
  \item Der Prozess wird korrekt weiter geführt, auch wenn der Server Abstürzen würde.
  \item Was passiert mit einer neuen Version eines Prozessen und so weiter.                      
\end{itemize}

Das sind alles relativ komplexe aufgaben, welche man im Normalfall nicht selber neu erfinden möchte.

Es folgt nun ein verglich verschiedener Engines. Die getroffene Auswahl ist weit weg von vollständig.



\subsection{activVOS}
Hersteller: Active Endpoints, Inc.

Sprache: java

Lizens: proprietär (setzt aber auf offenen standards)


Preise?

Dokumentation: Diverse Videos mit Demos und Tutorials.

Bemerkungen: Ist als eclipse Plugin aufgebaut. \\
Hat zum Beispiel ein Report designer. \\
Viele GUI's


\subsection{biztalk}
Hersteller: Microsoft

Sprache: .NET

Lizens: proprietär


Preise?

Dokumentation:

Bemerkungen:
Zusammenspiel mit Sharpoint.
nur für BPEL


\subsection{Apache ODE}
Hersteller: apache fundation

Sprache: java

Lizens: open source (Apache License 2.0 )

Bemerkungen:
Nur BPEL engine



\subsection{activiti}
Hersteller:

Sprache: java

Lizens: opensource (Apache License 2.0 )


Bemerkungen:
Nur BPMN engine
Schlank und aufgeräumt. 
Relativ neu aber Fork aus jBPM.


\subsection{ARIS}
Hersteller: software AG

Sprache: java

Lizens: proprietär

Preis?

Bemerkungen:
Unterstützt BPMN und BPEL. (xml nicht sichtbar/exportierbar...)
Viele GUI's. 
Monitoring, Reporting werkzeuge dabei.
Hauptsächlich für SAP entwickelt
Wird eventuel von <Firma> eingekauft/verwendet.
Programmierung und Konfiguration muss fast eingekauft werden.



\subsection{Zusammenfassung}
BPMN ist auf Grund der Business nähe eher für Laborprozesse geeignet als das BPEL.

Betrachtet man jetzt nur die BPMN Engines, bleibt von den oben genannten nur noch das ActiVOS, activit und ARIS übrig.

ActiVOS bietet vor allem zusätzliche GUI's an um die Konfiguration "einfacher" zu machen. Es geschieht viel im Hintergrund. 
So gesehen bietet es einen relativen geringen Mehrwert. 

ARIS ist das grosse rundum-sorglos-Paket. Leider wird einem auch alles Code relevante versteckt. Es daher nötig alles bei der anbietenden
Firma einzukaufen oder extra ein paar Schulungen beim Anbieter zu machen.

activit verzichtet auf eigene Gui's. Es bietet eine durchdachte Programmierschnitstelle an. Dadurch bleibt es schlank und beschränkt sich
auf das wesentliche, nämlich dem verwalten der Prozesse. Die Einarbeitung ist eher schwieriger, da direkt alles im Code implementiert wird.
Da es aber Opensource ist, gibt es bereits gute Beispiele wie eine Webapplikatione auszusehen hat.

Der monolytische Ansatz von grossen Systemen die alles können wollen, hat sich in der Vergangenheit nicht bewährt. Mit dieser Anforderung
sticht das activiti sehr positiv heraus. Weiter Tools die benötigt würden wie das Monitoring oder Reporting, können dann von anderen Tools
welche in dem Gebiet spezialisiert sind implementiert werden.
Nimmt man zudem Vaadin als Webframework, spielt es sehr gut mit activiti zusammen.


\subsection{Kompatibilität}
Trotz des BPMN Standards, haben die verschiedenen implementierunge der Engines eigene Schlüsselwörter, wie z.B für Formulare. 
Die BPMN Beschreibungen können daher nicht "einfach so" von einer zu anderen engine kopiert werden. 
Die Wahl der Engine ist demnach schon bindend.

Da aber immerhin die Geschäftslogik schön abgekapselt ist, ist das portieren einfacher, als wenn man alles in reinem Java programmiert
hätte.


\section{Beispiel Applikation}
Der Prozess beschreibt eine Jährliche Inventarkontrolle von Rückstellmustern.

Liste der Rückstellmuster wird aus dem übergeordneten System geholt (z.B LIMS). 

Liste wird anschliessend abgearbeitet. Dies kann von verschiedenne Personen Blockweise geschehen. Abarbeitung muss auch Offline passieren
können. Die Einzelnen Musterchecks erzeugen anschliessend einen Report zur Bestätigung und Nachverfolgbarkeit.

Zusammenfassende Liste mit Resultaten und eventuellen fails wird als Report gedruckt.

Reports werden in der elektronischen Datenbank abgelegt (anderes system). 

Dier Prozess wird von nun an als Beispiel herangezogen.

\subsection{Technologie Wahl}


\subsection{Webapplikation}

\subsection{bpmn}


\subsection{}



\section{Datennormalisierung}





\addcontentsline{toc}{section}{Literaturverzeichnis}
\begin{thebibliography}{99}

\bibitem{activiti} \url{http://www.activiti.org/}

\bibitem{omg} \url{http://www.omg.org/}




\end{thebibliography}

\end{document}

